\textbf{Piecewise linear regression for yield curve fitting}

Given a set of \emph{N} bonds tradable in market :

\(\left\{ T_{n},b_{n},a,c_{n},d_{n} \right\}\) for all \emph{n} belongs
to \emph{\{0,N-1\}}, where

\emph{T\textsubscript{n}} is bond maturity

\emph{b\textsubscript{n}} is current best bid of bond \emph{n}

\emph{a\textsubscript{n}} is current best ask of bond \emph{n}

\emph{c\textsubscript{n}} is previous close price of bond \emph{n}

\emph{d\textsubscript{n}} is dv01 of bond \emph{n}

As the objective is to fit the delta-yield curve
\(\mathrm{\Delta}y\left( \tau|{\mathrm{\Delta}y}_{\tau 0},{\mathrm{\Delta}y}_{\tau 1},\ldots{\mathrm{\Delta}y}_{\tau(M - 1)} \right)\),
with \emph{τ} being any tenor point on the delta-yield curve. The curve
is parameterized by delta-yield at a set of predefined time points, that
is \{\({\mathrm{\Delta}y}_{\tau m}\)\} where \emph{m} belongs to
\emph{\{0,M-1\}}. where delta-yield means the difference between current
yield from previous closing yield at a particular tenor point. This is a
piecewise linear model with \emph{M} control points, usually used
control points are :

\begin{quote}
\(\tau_{0} = 0\)

\(\tau_{1} = 1M\)

\(\tau_{2} = 3M\)

\(\tau_{3} = 6M\)

\(\tau_{4} = 1Y\)

\(\tau_{5} = 2Y\) \ldots{} etc
\end{quote}

The first step involves converting observed bond prices to delta-yields
using method like bisection or Newton Raphson (by solving
Internal-Rate-Return equation). Alternatively, we can approximate price
and yield relation using linear term of Taylor series, as long as the
pivot is close enough. Here we take the previous day closing price as
the pivot, for bond \emph{n} :

\(\frac{1}{2}(b_{n} + a_{n}) = c_{n} + d_{n}\mathrm{\Delta}y_{Tn}\)

Please note that \(\mathrm{\Delta}y\) is the difference between current
yield with closing yield of a particular bond, and different bonds have
different closing yield for previous day. We are actually performing a
regression (linear interpolation) in delta-yield space, rather than in
absolute-yield space (explanation in remark). Here are the notations :

\({\{\mathrm{\Delta}y}_{Tn}\}\) \emph{n = {[}0,N-1{]}} are observations
/ data points

\{\({\mathrm{\Delta}y}_{\tau m}\)\} \emph{m = {[}0,M-1{]}} are
parameters / control points

\(\mathrm{\Delta}y\left( \tau|\ldots \right)\) is the regression of
delta-yield curve at τ

The primary objective is to minimize weighted sum of price-error square
wr.t. control points :

\(\min_{\{\mathrm{\Delta}y_{\tau m}\}}\sum_{n = 0}^{N - 1}{w_{n}\left( c_{n} + d_{n}{\mathrm{\Delta}y}_{Tn} - \frac{1}{2}(b_{n} + a_{n}) \right)^{2}}\)

→
\(\min_{\left\{ {\mathrm{\Delta}y}_{\tau m} \right\}}\sum_{n = 0}^{N - 1}{w_{n}\left( c_{n} + d_{n}\left\lbrack {(1 - µ}_{Tn}){\mathrm{\Delta}y}_{\tau m} + µ_{Tn}{\mathrm{\Delta}y}_{\tau(m + 1)} \right\rbrack - \frac{1}{2}(b_{n} + a_{n}) \right)}^{2}\)
\emph{(1)}

Equation \emph{(1)} makes use of linear interpolation in piecewise
linear model, where maturity of bond \emph{n} lies between 2 tenor
points in delta-yield curve \(\tau_{m} \leq T_{n} \leq \tau_{m + 1},\)
where \emph{w\textsubscript{n}} is the weight of the
n\textsuperscript{th} data point (an example can be inverse of its
market spread), while \emph{µ\textsubscript{Tn}} is for linear
interpolation weight.

\begin{quote}
\(w_{n} = \left( \frac{1}{a_{n} - b_{n}} \right)^{2}\)
\(µ_{Tn} = \frac{T_{n} - \tau_{m}}{\tau_{m + 1} - \tau_{m}}\)
\end{quote}

Rewriting \emph{(1)} as matrix form :

\(\min_{Y}(A\mathrm{\Delta}Y - B)^{T}W(A\mathrm{\Delta}Y - B)\)
\emph{where W is diagonal matrix of weight} \emph{(2)}

where \emph{A} is \emph{N×M} design matrix, \(\mathrm{\Delta}\)\emph{Y}
is \emph{M} parameter vector, \emph{B} is \emph{N} output vector.

\begin{quote}
\[A = \begin{bmatrix}
0 & d_{1}\frac{\tau_{2} - T_{1}}{\tau_{2} - \tau_{1}} & d_{1}\frac{T_{1} - \tau_{1}}{\tau_{2} - \tau_{1}} & 0 & \ldots & 0 & 0 \\
0 & 0 & d_{2}\frac{\tau_{3} - T_{2}}{\tau_{3} - \tau_{2}} & d_{2}\frac{{T_{2} - \tau}_{2}}{\tau_{3} - \tau_{2}} & \ldots & 0 & 0 \\
d_{3}\frac{\tau_{1} - T_{3}}{\tau_{1} - \tau_{0}} & d_{3}\frac{{T_{3} - \tau}_{0}}{\tau_{1} - \tau_{0}} & 0 & 0 & \ldots & 0 & 0 \\
\ldots & \ldots & \ldots & \ldots & \ldots & \ldots & \ldots \\
0 & 0 & 0 & 0 & \ldots & d_{N}\frac{\tau_{M + 1} - T_{N}}{\tau_{M + 1} - \tau_{M}} & d_{N}\frac{T_{N} - \tau_{M}}{\tau_{M + 1} - \tau_{M}}
\end{bmatrix}\]

\(\mathrm{\Delta}Y = \begin{bmatrix}
{\mathrm{\Delta}y}_{\tau 0} \\
{\mathrm{\Delta}y}_{\tau 1} \\
\mathrm{\Delta}y_{\tau 2} \\
\ldots \\
{\mathrm{\Delta}y}_{\tau M}
\end{bmatrix}\) B\(= \begin{bmatrix}
\frac{1}{2}\left( b_{1} + a_{1} \right) - c_{1} \\
\frac{1}{2}\left( b_{2} + a_{2} \right) - c_{2} \\
\frac{1}{2}\left( b_{3} + a_{3} \right) - c_{3} \\
\ldots \\
\frac{1}{2}\left( b_{N} + a_{N} \right) - c_{N}
\end{bmatrix}\)
\end{quote}

In order to enhance regression, we apply linear transformation \emph{T}
(an orthonormal matrix) on parameter space, transforming it from
delta-yield \(\mathrm{\Delta}Y\) to \(T\mathrm{\Delta}Y\). \emph{T} is
actually \emph{PCA} rotation of the \emph{M} dimensional delta-yield
space. In other words, instead of doing regression in delta-yield, we do
it in risk factor space \(\mathrm{\Delta}Y'\), spanned by yield-curve
level shift, tilt, curvature together with other higher order factors.

\(\min_{\mathrm{\Delta}Y}(A\mathrm{\Delta}Y - B)^{T}W(A\mathrm{\Delta}Y - B)\)

=
\(\min_{\mathrm{\Delta}Y}\left( AT^{- 1}T\mathrm{\Delta}Y - B \right)^{T}W(AT^{- 1}T\mathrm{\Delta}Y - B)\)

=
\(\min_{\mathrm{\Delta}Y'}\left( A^{'}{\mathrm{\Delta}Y}^{'} - B \right)^{T}W(A'\mathrm{\Delta}Y' - B)\)
\emph{(3)}

\emph{where} \(A^{'} = AT^{- 1}\) \emph{= interpolation weight}

\({\mathrm{\Delta}Y}^{'} = T\mathrm{\Delta}Y\) \emph{= values on grid
points}

\(B\) \emph{= values on observation points}

Furthermore, \emph{T} transformation usually works with penalty
\emph{P}, a \emph{M} by \emph{M} diagonal matrix, with diagonal values
\emph{{[}0, 0, 0, 1, 1, \ldots, 2, 2, \ldots{]}}, which means no penalty
for the first 3 risk factors : curve level shift, curve tilt and curve
curvature, minor penalties for next few risk factors, higher penalties
for high order risk factors. As we can see from the following objective
function, higher penalty implies stronger force to make corresponding
element in \(\mathrm{\Delta}Y^{'} = T\mathrm{\Delta}Y\) to zero, i.e.
looking for a smoother yield curve. Objective becomes minimization of
\emph{J(}\(\mathrm{\Delta}\)\emph{Y')} :

\(J(\mathrm{\Delta}Y') = \left( A^{'}{\mathrm{\Delta}Y}^{'} - B \right)^{T}W\left( A^{'}\mathrm{\Delta}Y^{'} - B \right) + {{\mathrm{\Delta}Y}^{'}}^{T}P\mathrm{\Delta}Y^{'}\)
\emph{(4)}

\(\frac{dJ}{d\mathrm{\Delta}Y'} = {2A^{'}}^{T}W\left( A^{'}\mathrm{\Delta}{Y'}_{opt} - B \right) + 2P\mathrm{\Delta}{Y'}_{opt} = 0\)

\({A^{'}}^{T}WA^{'}{\mathrm{\Delta}Y'}_{opt} + P{\mathrm{\Delta}Y'}_{opt} - {A^{'}}^{T}WB = 0\)

\({\mathrm{\Delta}Y'}_{opt} = \left( {A^{'}}^{T}WA^{'} + P \right)^{- 1}\left( {A^{'}}^{T}WB \right)\)

In practice, we \hl{do not need to solve for delta-yield value} (or even
absolute yield rate) in the yield curve, our objective is, given a noisy
observation of the mid price of a bond, we can remove noise by adjusting
to the fitting result (i.e. averaging out the noise of observations of
all bonds).

Adjusted mid price \emph{m\textsubscript{n,adjusted}} of \hl{bond
\emph{n} in the training set} is found by :

\({A^{'}}_{n}{{\mathrm{\Delta}Y}^{'}}_{opt} = B_{n} = m_{n,adjusted} - c_{n}\)

\(m_{n,adjusted} = {A'}_{n}{\mathrm{\Delta}Y'}_{opt} + c_{n}\)

\emph{where} \({A^{'}}_{n} = \lbrack\begin{matrix}
0 & d_{n}\frac{{\tau_{2} - T}_{n}}{\tau_{2} - \tau_{1}} & d_{n}\frac{T_{n} - \tau_{1}}{\tau_{2} - \tau_{1}} & 0 & \ldots & 0
\end{matrix}\rbrack\)

\emph{assuming} \(\tau_{1} \leq T_{n} \leq \tau_{2}\)

Adjusted mid price \emph{m\textsubscript{adjusted}} of \hl{any bond
\emph{(with maturity T and close price c)}} is found by :

\(m_{adjusted} = {A'}_{1row}{\mathrm{\Delta}Y'}_{opt} + c\)

\emph{where} \({A^{'}}_{1row} = \lbrack\begin{matrix}
0 & d_{n}\frac{\tau_{2} - T}{\tau_{2} - \tau_{1}} & d_{n}\frac{T - \tau_{1}}{\tau_{2} - \tau_{1}} & 0 & \ldots & 0
\end{matrix}\rbrack\)

\emph{assuming} \(\tau_{1} \leq T \leq \tau_{2}\)

\textbf{Implementation}

We can divide the bond price adjustment into 2 steps :

\({\mathrm{\Delta}Y'}_{opt} = \left( {A^{'}}^{T}WA^{'} + P \right)^{- 1}\left( {A^{'}}^{T}WB \right)\)
\emph{training}

\(m_{adjusted} = {A'}_{1row}{\mathrm{\Delta}Y'}_{opt} + c\)
\emph{testing}

We re-arrange these 2 steps into the following for latency :

• slow but very rare step, invoked when there is change in \emph{A, W,
P} (not likely to happen)

• fast but frequent step, invoked when there is change in market data
\emph{B} (happen every tick)

\(R = \left( {A^{'}}^{T}WA^{'} + P \right)^{- 1}\left( {A^{'}}^{T}\sqrt{W} \right)\)
\emph{slow but rare : pseudo inverse}

\(m_{adjusted} = {A^{'}}_{1row}(R\sqrt{W}B) + c\) \emph{fast but freq :
per tick update}

\textbf{\hfill\break
}

\textbf{Correspondence in the code\\
}Notation above uses row major matrix :

\emph{N = number of data = number of rows in A}

\emph{M = dimension of parameter = number of columns in A = sizeof Y}

\emph{A = design matrix with row data}

However the convention in ffwk code is transpose of my convention.
Please convert them appropriately when reading code. In the following
table, first column is variable in code, second column is variable in my
equations above.

In function FSimpleLinearCurveMid::initialize()

m\_wmat \(T^{- 1}\) \emph{PCA transformation}

m\_trans \(A^{'} = AT^{- 1}\) \emph{PCA transformed \textbf{design
matrix}}

m\_trans2 \({A^{'}}^{T}A^{'} + P\) \emph{\textbf{Gram matrix} for LS}

m\_transinv \(\left( {A^{'}}^{T}A^{'} + P \right)^{- 1}\)
\emph{\textbf{inverse Gram matrix} for LS}

In function FSimpleLinearCurveMid::updateRotation()

m\_midwght \(\sqrt{W}\) \emph{diagonal weight matrix i.e.} \(W = W^{T}\)

m\_trans\_scaled \({\sqrt{W}A}^{'}\) \emph{PCA transformed
\textbf{design matrix}}

m\_trans2 \({A^{'}}^{T}{WA}^{'} + P\) \emph{\textbf{Gram matrix} for LS}

m\_transinv \(\left( {A^{'}}^{T}WA^{'} + P \right)^{- 1}\)
\emph{\textbf{inverse Gram matrix} for LS}

m\_rotation
\(\left( {A^{'}}^{T}WA^{'} + P \right)^{- 1}{A^{'}}^{T}\sqrt{W}\)
\emph{\textbf{pseudo inverse} for LS}

In function FSimpleLinearCurveMid::operator()

m\_res \(R\sqrt{W}B\) \emph{optimized parameters}

m\_tmp \(A^{'}(R\sqrt{W}B)\)

value() \(A^{'}(R\sqrt{W}B) + c\) \emph{adjusted price to fitted line}

In function FSimpleLinearCurveMid::initialize(), it finds :

• \(\left( {A^{'}}^{T}A^{'} + P \right)^{- 1}\) with all-1-weight once,
and then

• \(\left( {A^{'}}^{T}WA^{'} + P \right)^{- 1}\) once more via calling
updateRotation()

Why is that? Are there redundancy? Slightly, the objective of first call
is not for solving inverse, its main aim is to get a determinant value
under all-1-weight condition, this determinant is saved and used for
checking if ill-condition happens when weight W is changing in the
lifetime of the functor.

\textbf{\hfill\break
}

\textbf{Parameter space and Relative value (RV)}

Intuitively, we are supposed do regression in yield curve space \(Y\),
but we convert it to delta yield space \(\mathrm{\Delta}Y\). Then we
further apply a PCA rotation to risk factor face
\(\mathrm{\Delta}Y^{'} = T\mathrm{\Delta}Y\). We eventually optimize in
\(\mathrm{\Delta}Y^{'}\). Why is that? There are 2 conversion steps :

• conversion from yield to delta yield

• conversion from delta yield to risk factor (PCA)

The reason for PCA is obvious. This is because we want to model main
yield curve factor, rather than modelling the yield value on the nodes.
Yield curve factor like level shift, tilt and curvature have their
economic meaning, like market expectation of economics, inflation etc.

The reason for conversion of yield to delta yield is because of relative
value (RV) of bonds. Theoretically all bonds should have price following
the same yield curve. But this is not the case in reality. There is a
spread between the market traded price of a bond with the present value
implied from the yield curve, i.e. different bonds with same maturity,
same or different coupons, may have different implied yields, because of
various reasons (these reasons are called \hl{relative values}, or
\hl{basis} in ffwk model). If a bond is traded at a higher price than
its peer in the yield curve, it is said to trade rich (at premium).
Otherwise it is said to trade cheap.

Relative values include :

• country basis : sovereign yield differential at the same maturity

• on the run bond vs off the run bond : former has a premium

• cheapest to deliver (CTD) bond has a premium

• repo specialized bond has a premium

• asset swap spread (how?)

Therefore the market price of a bond is the price given by yield curve
plus sum of the above relative values (or sum of above basis).

\emph{bond price = yield curve price + country basis + on the run
premium +}

\begin{quote}
\emph{CTD premium + repo premium}

= \(A^{'}(R\sqrt{W}B) + c\) \emph{+ country basis + on the run premium
+}

\emph{CTD premium + repo premium}
\end{quote}

Remarks

• some firms speculates these relative values, while keeping exposure
neutral to yield curve risk

• alpha price model in ffwk are actually modelling the above basis
(please see json in model)
